%=====================================================================
\chapter{Results Discussion and Research Significance}
%=====================================================================

\section{Consolidated Outcomes Against the Objectives}

Table~\ref{tab:outcomes} summarises what each objective produced and the status of the
evidence behind it. The final column is stated deliberately, because the strength of a
result and the confidence that can be placed in it are separate properties and are treated
separately throughout this work.

\begin{table}[!htb]
\centering
\caption[Objective-wise consolidated outcomes and evidence status]{Objective-wise consolidated outcomes and the status of the supporting evidence}
\label{tab:outcomes}
\tabfont
\begin{tabular}{@{}L{0.85cm}L{4.4cm}L{4.4cm}L{4.5cm}@{}}
\toprule
\textbf{Obj.} & \textbf{Principal outcome} & \textbf{Key quantitative evidence} &
\textbf{Status and limits} \\
\midrule
RO1 & Four-stream fusion model with adaptive, then uncertainty-conditioned, source
weighting & MAE 0.018, $R^2$ 0.988, DA 84.7\% against 76.3\% for a fixed-weight ensemble
on identical inputs; both textual sources contribute 10.9\% of $R^2$ superadditively &
Established. Index-level, one-day horizon; macro interpolation understates announcement-day
effects \\
RO2 & Hybrid LLM sentiment and indicator framework in a reproducible environment &
Sentiment is the highest-gain feature; correlation with price 0.04, with momentum 0.74;
accuracy 47.72\% against a 46.43\% baseline mean; highest prediction confidence &
Established but bounded. All models near the random classifier at a single-name daily
horizon \\
RO3 & Causal, multi-granularity transparency and a measured account of sentiment noise &
Causal graph density 0.850 with an identifiable hub; correct external sign in 16 of 16
series-by-period tests; market states separated at $p=1.8\times10^{-8}$ & Established for
construct validity. Predictive content of the composite is separately measured and is not
significant \\
RO4 & Regime-conditional decision framework, ablated at matched risk and tested across
markets & Sharpe 0.878 against 0.628; drawdown $-18.98\%$ against $-36.08\%$; wins in
18/18 specifications, 18/19 years, 4/4 stress episodes; holdout Sharpe 1.614 & Established
as consistent, not as statistically significant. Bootstrap interval
$[-0.102, 0.574]$ includes zero \\
\bottomrule
\end{tabular}
\end{table}

The final column separates what the research demonstrates from what it deliberately
does not claim. Read as a ledger the pattern is consistent: adaptive fusion beats
fixed-weight combination, but index-level accuracy does not transfer to single-name
daily direction; language-model sentiment is orthogonal and highest-importance, but
does not overturn market efficiency at a daily horizon; the sentiment construct is
validated, but a validated construct does not thereby forecast returns; regime risk
timing improves both Sharpe ratio and drawdown at matched exposure, but the Sharpe
gain is not statistically significant on eighteen years of data; and the defensive
behaviour transfers to a second universe, though the framework does not add value in
a concentrated single-factor one. These bounding statements are not shortcomings but
the boundary of the contribution, and stating them is what makes the positive claims
credible.

\section{Reconciling the Divergent Performance Across Studies}

A reader comparing the objectives will notice an apparent contradiction: Objective~1
reports 84.7 per cent directional accuracy while Objective~2 reports 47.72 per cent. The
two numbers are not in conflict, and understanding why is central to the significance of
this research.

They answer different questions on different objects. Objective~1 predicts the next-day
level of a diversified index from four fused sources; an index aggregates hundreds of
constituents, so idiosyncratic noise largely cancels and the remaining signal is dominated
by persistent market-wide state, which fusion captures well. Objective~2 predicts the
next-day direction of an individual equity, which carries its full idiosyncratic
component, has close to balanced classes and offers no aggregation to average the noise
away. Predictability is therefore expected to be far lower, and this is what the empirical
asset pricing literature reports~\cite{gu2020}.

Three implications follow, and together they define the trajectory of the research
programme. First, an accuracy figure is meaningless without its object and horizon.
Second, the level at which a system is applied should match where the signal actually
lives: aggregate state is predictable, single-name daily direction is largely not. Third,
and most consequentially, if the signal channel is weak at the level where decisions are
taken, value must be created at the decision layer rather than the prediction layer.
Objective~4 tests exactly that, and the matched-exposure ablation confirms it: risk timing
works where signal selection does not. The apparent contradiction between the two accuracy
figures is therefore not an inconsistency but the finding that organises the work.

\section{Research Significance}

\textbf{Methodological significance.} The work advances multi-source fusion from a
question of what to combine to a question of how to combine under changing reliability.
Expressing fusion weights as a function of measured, time-varying predictive
uncertainty~\cite{farimani2024} gives a source-weighting rule that self-corrects within
days rather than at the next retraining cycle, and distribution-aware normalisation
prevents a heavy-tailed input from dominating the fused representation. Separating the
allocation decision into relative composition and total risk is a second contribution: it
makes each layer independently falsifiable, and without it the central ablation of this
thesis could not have been constructed at all.

\textbf{Evaluation significance.} Three practices used here are uncommon in this
literature and are the most transferable part of the work. The first is ablation at
matched risk, which separates the timing of risk from its level; comparing a de-risked
strategy against a fully invested benchmark confounds the two and credits a framework for
merely holding less. The second is disclosure of the entire specification sweep rather
than its best cell, the minimum required once multiple testing is taken
seriously~\cite{harvey2015,valls2026}. The third is separation of construct validity from
predictive usefulness: the sentiment composite is shown to measure sentiment and,
separately, shown not to forecast returns. Conflating those questions is a recurring
error, and keeping them apart changes what a null result means.

\textbf{Engineering significance.} Two failure modes were diagnosed that generalise beyond
this framework. The first is scale mismatch in signal-to-allocation conversion: passing an
unscaled indicator into a mean-variance objective lets the ratio of units between the
linear and quadratic terms, rather than the information in the signal, determine the
solution, producing a degenerate single-asset portfolio; information-ratio
scaling~\cite{grinold2000} removes it. The second is an inert risk layer: an absolute
volatility target calibrated for a concentrated equity portfolio binds in one rebalance
out of 224 on a diversified universe, so the layer is present, consumes computation and
does nothing, whereas a relative budget is scale-free and remains active. Mixture-model
label anchoring has the same character: omitting it produces no error message, only a
silently scrambled regime series.

\textbf{Practical significance.} For a practitioner the results support a specific and
limited prescription. Fuse sources, but weight them by measured reliability rather than by
historical average importance. Use language-model sentiment for what it demonstrably
provides --- orthogonal information and better-calibrated confidence --- rather than for a
directional edge it does not provide at a daily horizon. Prefer causally selected features
where the structure is recoverable, because they survive regime change better and can be
audited. Implement a regime risk budget in relative rather than absolute form, evaluate it
against a volatility-matched benchmark, and expect it to help only where the universe
contains genuine dispersion in asset volatility. For the Indian market specifically, the
measured weights matter: macroeconomic indicators carry 25 per cent of average attention,
reflecting the influence of policy rates, inflation and industrial output on valuations,
sentiment sources jointly carry 40 per cent with the social media share rising sharply at
market extremes, and sectoral predictability varies systematically with information
availability and factor exposure. These are measured properties of that information
environment rather than assumptions imported from elsewhere.

%=====================================================================
\chapter{Conclusion and Future Scope}
%=====================================================================

\section{Conclusion}

This research set out to integrate multi-source data with sentiment analysis and language
models in order to improve stock market decision making, and to determine which parts of
such an integration actually contribute to the outcome. Both aims are met, though not
entirely in the direction originally anticipated, and the divergence is itself the most
useful result.

A four-stream fusion model integrating financial metrics, news sentiment, social media
sentiment and macroeconomic indicators was developed and evaluated on ten years of Indian
market data across ten sectoral indices. It attains a normalised mean absolute error of
0.018, a coefficient of determination of 0.988 and a directional accuracy of 84.7 per cent
against 76.3 per cent for a fixed-weight ensemble over identical inputs, with stability
confirmed by walk-forward evaluation at $0.0016 \pm 0.0004$ normalised error. Removing
both textual sources costs 10.9 per cent of explained variance, exceeding the sum of their
individual removals --- the quantitative signature of complementary information. The
fusion rule was then advanced from learned average relevance to measured instantaneous
reliability through a Bayesian formulation in which weights are obtained by a softmax over
negative predictive uncertainties.

A hybrid framework coupling a transformer language model with fifteen conventional
technical indicators was implemented in a reproducible computational environment and
evaluated under a strict temporal protocol. Language-model sentiment emerged as the
highest-gain feature in the fused space, correlating 0.74 with momentum yet only 0.04 with
price, establishing that it supplies orthogonal information. The gain in raw accuracy is
modest --- 47.72 per cent against a baseline mean of 46.43 per cent --- and all models
operate close to the random classifier at a single-name daily horizon. That limit is
reported rather than concealed, and it redirected the research toward the decision layer.

Transparency was established structurally. A directed causal adjacency matrix recovered by
conditional independence testing yields a graph of density 0.850 with an identifiable
causal hub, and demonstrates that strong directed influence can coexist with weak linear
correlation --- the precise justification for causal rather than correlational feature
selection. Sentiment noise was quantified rather than assumed: the composite carries the
expected sign against four independently published series in all sixteen
series-by-sub-period combinations and separates the identified market states at
$p = 1.8 \times 10^{-8}$, while its predictive content, measured separately, is not
significant. Keeping construct validity and forecasting power apart is a contribution in
its own right.

The decision framework was evaluated over eighteen years on ten liquid multi-asset
instruments under a walk-forward protocol with transaction costs. It attains a Sharpe
ratio of 0.878 against 0.628 for an equal-weighted benchmark and reduces maximum drawdown
from $-36.1$ to $-19.0$ per cent, outperforming in all four stress episodes, in 18 of 19
calendar years on drawdown and in all 18 specifications examined; on a 665-session holdout
that postdates every design decision it records a higher Sharpe ratio and a smaller
drawdown than all three benchmarks. Ablation at matched average exposure locates the
source precisely: timing risk by identified market state improves both the Sharpe ratio
and the drawdown with no significant change in mean return, whereas the signal fusion
layer contributes nothing detectable. Cross-market transfer reproduces the defensive
behaviour in a concentrated mega-capitalisation universe while losing on risk-adjusted
return, bounding applicability to universes with genuine dispersion in asset volatility.

The overall conclusion is precise rather than expansive. Multi-source integration and
language-model sentiment demonstrably improve the description of the market, and causal
structure demonstrably improves the transparency of that description. At the horizon and
granularity at which decisions are actually taken, however, the measurable value of the
system arises at the allocation layer, through regime-conditional risk timing, rather than
at the signal layer. A framework whose every layer has been separately tested, and whose
applicability has been bounded by evidence, is a more useful engineering artefact than one
whose aggregate performance conceals which of its components is doing the work.

\section{Future Scope}

\begin{itemize}
\item \textbf{Textual sentiment in a universe with instrument-level news coverage.} The
null reported for the market-derived composite is a null about that proxy, not about
textual sentiment. Applying the fusion hypothesis where a time-stamped, instrument-level
news corpus exists would convert a proxy null into a genuine test, and would connect the
language-model pipeline of Objective~2 directly to the allocation layer of
Objective~4~\cite{officioso2026,liu2024llmreview}.

\item \textbf{Bayesian neural networks for richer uncertainty.} The present uncertainty
estimate is a rolling error variance, which captures aleatoric variation. Replacing it
with proper Bayesian neural expert models would additionally expose epistemic uncertainty,
producing a confidence measure that distinguishes a noisy environment from an unfamiliar
one~\cite{khuwaja2023}.

\item \textbf{Non-stationary debiasing and scalable causal discovery.} The stationarity
assumption in the multi-scale debiasing layer fails at exactly the moments that matter, as
the 77.67 per cent peak volatility during the Union Budget window shows; an adaptive
formulation whose noise model is itself state-dependent is the direct remedy. In parallel,
extending conditional independence testing from five instruments to the full
fifty-constituent index requires an algorithm whose cost does not grow quadratically,
which is a prerequisite for causal transparency at production
scale~\cite{liu2026dynamic,spirtes2000}.

\item \textbf{Adaptation-based cross-market transfer.} The cross-market evidence reported
here comes from applying the framework to a second universe without modification, which
tests whether the mechanism generalises but forgoes any benefit from adapting to the target
market. The natural extension is genuine transfer learning: pre-training the regime model
and the source-reliability estimates on a data-rich market and fine-tuning them on a
thinner one, so that an emerging-market deployment inherits structure learned elsewhere
instead of starting from its own limited history.

\item \textbf{Learned risk budgets and explicit state persistence.} The regime multipliers
are fixed at 1.00, 0.75 and 0.50; learning them --- reinforcement learning being the
natural candidate given the sequential structure --- would adapt the budget to the
universe, subject to the overfitting controls this work argues are
mandatory~\cite{harvey2015}. Replacing the Gaussian mixture with a hidden Markov or jump
formulation would represent regime persistence directly rather than recovering it
incidentally~\cite{luo2026regime,hamilton1989}.

\item \textbf{Regional-language and higher-frequency sentiment.} Sentiment extraction is
currently English-only and daily. A large share of Indian retail participants read
financial news in regional languages, and intraday resolution would expose incorporation
dynamics that daily aggregation averages away.

\item \textbf{Deployment engineering.} Inference cost is negligible; the binding
constraints are the latency of retrieving and scoring text at scale, feed rate limits,
coordinated manipulation of social signals, transaction costs and compliance with
algorithmic trading regulation. These are treated as first-class engineering problems
rather than as afterthoughts.
\end{itemize}

%=====================================================================
\chapter{Organisation of the Thesis}
%=====================================================================

The thesis is organised into nine chapters.

\textbf{Chapter I --- Introduction.} Introduces stock market decision making as an
information-integration problem, describes the structure of the Indian equity market, and
states the motivation for combining quantitative, textual and macroeconomic evidence
within a single system.

\textbf{Chapter II --- Background and Preliminaries.} Presents the technical foundations:
technical indicator construction, sentiment extraction from lexicon methods through to
transformer language models, fusion strategies, regime modelling, and the evaluation
metrics and statistical procedures used throughout the thesis.

\textbf{Chapter III --- Literature Survey and Research Gaps.} Reviews multi-source fusion,
language-model sentiment analysis, explainable and causal financial modelling, and
decision-level allocation, and derives the five research gaps that define the scope of the
work.

\textbf{Chapter IV --- Multi-Source Data Integration Framework.} Develops the
Multi-Source Fusion Network: data acquisition and temporal alignment, modality-specific
encoders, cross-modal attention, the Bi-LSTM predictor, the experimental setup, the
sectoral and ablation results, and the walk-forward robustness analysis.

\textbf{Chapter V --- Dynamic Uncertainty-Weighted Fusion.} Develops the Bayesian
formulation of the fusion rule: expert models per source, time-varying uncertainty
estimation, softmax weighting over negative uncertainties, distribution-aware
normalisation, and comparison against static early, late and attention-based fusion
baselines.

\textbf{Chapter VI --- Hybrid Language Model and Indicator Framework.} Develops the
sentiment extraction pipeline, feature-level fusion with technical indicators, model
training and comparison, feature-importance and sentiment-regime analysis, confidence
calibration and the simulated trading evaluation.

\textbf{Chapter VII --- Explainable and Causal Framework.} Develops causal graph discovery,
multi-granularity temporal memory encoding and credibilistic optimisation; reports the
recovered causal structure, the correlation-versus-causality contrast, the sentiment noise
diagnostics and the sentiment-to-price temporal dynamics.

\textbf{Chapter VIII --- Decision-Theoretic Allocation and Cross-Market Validation.}
Develops regime identification, constrained portfolio optimisation and the regime risk
budget; reports the eighteen-year evaluation, the specification sweep, the
matched-exposure ablation, the stress-episode analysis, the holdout and the cross-market
transfer test.

\textbf{Chapter IX --- Conclusion and Future Scope.} Consolidates the findings across the
four objectives, states the limits of the evidence, and sets out the future research
directions.

%=====================================================================
\chapter*{Publications on PhD Work}
\addcontentsline{toc}{chapter}{Publications on PhD Work}
%=====================================================================

\textbf{International Journals}

\begin{enumerate}[label={[P\arabic*]},leftmargin=2.6em,start=3]
\item A. Pardeshi and S. Deshmukh, ``Fusing LLM-Extracted News Sentiment and Market Data
for Enhanced Stock Movement Prediction,'' \emph{Liberte Journal}, ISSN 0024-2020,
vol.~14, no.~6, pp.~424--439, 2026. \textbf{(Published)}\\
\indexing{\LiberteIndexing} \\ \link{\LiberteURL}

\item A. Pardeshi and S. Deshmukh, ``Dynamic Uncertainty-Weighted Fusion for Multi-Source
Financial Data: A Bayesian Framework for Time-Varying Source Reliability,''
\emph{International Journal of Computer Information Systems and Industrial Management
Applications} (IJCISIM). \textbf{(Communicated, \IJCISIMDate)}\\
\link{\IJCISIMURL}

\item A. Pardeshi and S. Deshmukh, ``From Signal Fusion to Asset Allocation: A
Decision-Theoretic Model for Portfolio Construction Under Regime-Based Sentiment and
Volatility,'' \emph{International Journal of Applied Decision Sciences} (Inderscience),
manuscript IJADS-312370. \textbf{(Under review --- revised manuscript submitted after a
major-revision decision)}\\
\link{\IJADSURL}
\end{enumerate}

\textbf{International Conferences}

\begin{enumerate}[label={[P\arabic*]},leftmargin=2.6em]
\item A. Pardeshi and S. Deshmukh, ``A Multi-Source Data Fusion Framework for Enhanced
Stock Market Prediction,'' \emph{2026 International Conference on Sustainability,
Innovation and Technology (ICSIT)}, \ICSITVenue, IEEE, 2026, \ICSITPages.\\
\link{\ICSITURL}

\item A. Pardeshi and S. Deshmukh, ``A Dynamic-Causal Hybrid Framework for NIFTY50 Stock
Decision Making: A Multi-Granularity Temporal Learning Approach,'' \DCHFVenue.
\textbf{(Communicated)}\\
\link{\DCHFURL}
\setcounter{enumi}{5}
\item A. Pardeshi and S. Deshmukh, ``Deep Learning in Stock Market Forecasting:
Comparative Insights and Future Directions,'' \emph{2025 International Conference on
Emerging Trends in Industry 4.0 Technologies (ICETI4T)}, Navi Mumbai, India, IEEE,
June 2025, \ICETIPages.\\
\link{\ICETIURL}
\end{enumerate}

\ifHasIPR
\textbf{Copyrights and Patents}

\begin{itemize}
\item \IPRentries
\end{itemize}
\fi

\vspace{0.4em}
\textit{Note:} The labels \pubcite{P1}--\pubcite{P6} are used throughout this report to
refer to the above works. Publication status is as of the date of submission of this
report.

%=====================================================================
\clearpage
\begin{footnotesize}
\setstretch{1.0}
\setlength{\parskip}{0pt}
\setlength{\columnsep}{16pt}
\begin{multicols}{2}
\begin{thebibliography}{99}
\setlength{\itemsep}{1.5pt}\setlength{\parsep}{0pt}
\addcontentsline{toc}{chapter}{References}

\bibitem{snasel2024} V. Sn\'a\v{s}el, J. D. Vel\'asquez, M. Pant, D. Georgiou and L. Kong,
``A generalization of multi-source fusion-based framework to stock selection,''
\emph{Information Fusion}, vol.~102, p.~102018, 2024.

\bibitem{long2024} W. Long, J. Gao, K. Bai and Z. Lu, ``A hybrid model for stock price
prediction based on multi-view heterogeneous data,'' \emph{Financial Innovation}, 2024.

\bibitem{zhang2022fusion} C. Zhang, N. N. A. Sjarif and R. B. Ibrahim, ``Decision fusion
for stock market prediction: a systematic review,'' \emph{IEEE Access}, vol.~10,
pp.~81364--81379, 2022.

\bibitem{farimani2024} S. Anbaee Farimani, M. V. Jahan and A. Milani Fard, ``An adaptive
multimodal learning model for financial market price prediction,'' \emph{IEEE Access},
vol.~12, pp.~121846--121863, 2024.

\bibitem{moodi2024} F. Moodi, A. J. Rafsanjani, S. Zarifzadeh and M. A. Zare Chahooki,
``Fusion of technical indicators and sentiment analysis in a hybrid framework of deep
learning models for stock price movement prediction,'' \emph{IEEE Access}, vol.~12,
pp.~195696--195709, 2024.

\bibitem{liu2019attention} G. Liu and X. Wang, ``A numerical-based attention method for
stock market prediction with dual information,'' \emph{IEEE Access}, vol.~7,
pp.~7357--7367, 2019.

\bibitem{mu2023investor} G. Mu, N. Gao, Y. Wang and L. Dai, ``A stock price prediction
model based on investor sentiment and optimized deep learning,'' \emph{IEEE Access},
vol.~11, pp.~51353--51367, 2023.

\bibitem{haryono2023} A. T. Haryono, R. Sarno and K. R. Sungkono, ``Transformer-gated
recurrent unit method for predicting stock price based on news sentiments and technical
indicators,'' \emph{IEEE Access}, vol.~11, pp.~77132--77146, 2023.

\bibitem{haryono2024} A. T. Haryono, R. Sarno, R. N. E. Anggraini and K. R. Sungkono,
``Permuted temporal Kolmogorov--Arnold networks for stock price forecasting using
generative aspect-based sentiment analysis,'' \emph{IEEE Access}, vol.~12,
pp.~178672--178689, 2024.

\bibitem{tai2022} W. Tai, T. Zhong, Y. Mo and F. Zhou, ``Learning sentimental and
financial signals with normalizing flows for stock movement prediction,'' \emph{IEEE
Signal Processing Letters}, vol.~29, pp.~414--418, 2022.

\bibitem{bacco2024} L. Bacco, L. Petrosino, D. Arganese, L. Vollero, M. Papi and
M. Merone, ``Investigating stock prediction using LSTM networks and sentiment analysis of
tweets under high uncertainty,'' \emph{IEEE Access}, vol.~12, pp.~122239--122248, 2024.

\bibitem{jung2024} H. D. Jung and B. Jang, ``Enhancing financial sentiment analysis
ability of language model via targeted numerical change-related masking,'' \emph{IEEE
Access}, vol.~12, pp.~50809--50820, 2024.

\bibitem{liu2024llmreview} C. Liu, A. Arulappan, R. Naha, A. Mahanti, J. Kamruzzaman and
I.-H. Ra, ``Large language models and sentiment analysis in financial markets: a review,
datasets, and case study,'' \emph{IEEE Access}, vol.~12, pp.~134041--134061, 2024.

\bibitem{onozo2024} L. R. \'On\'oz\'o, F. V. Arthur and B. Gyires-T\'oth, ``Leveraging
LLMs for financial news analysis and macroeconomic indicator nowcasting,'' \emph{IEEE
Access}, vol.~12, pp.~160529--160547, 2024.

\bibitem{khuwaja2023} P. Khuwaja, S. A. Khowaja and K. Dev, ``Adversarial learning
networks for FinTech applications using heterogeneous data sources,'' \emph{IEEE Internet
of Things Journal}, vol.~10, no.~3, pp.~2194--2201, 2023.

\bibitem{vargas2022} G. Vargas, L. Silvestre, L. Rigo J\'unior and H. Rocha, ``B3 stock
price prediction using LSTM neural networks and sentiment analysis,'' \emph{IEEE Latin
America Transactions}, vol.~20, no.~7, pp.~1067--1074, 2022.

\bibitem{wang2024spcm} J. Wang and Z. Chen, ``SPCM: a machine learning approach for
sentiment-based stock recommendation system,'' \emph{IEEE Access}, vol.~12,
pp.~14116--14129, 2024.

\bibitem{shang2025} L. Shang, ``Sentiment-driven Bitcoin price range forecasting:
enhancing CART decision trees with high-dimensional indicators and Twitter dynamics,''
\emph{IEEE Access}, vol.~13, pp.~60508--60518, 2025.

\bibitem{sami2025} H. M. Sami, M. Shah Jalal, M. F. Kabir and M. N. Huda, ``IntPort: an
intelligent portfolio construction technique based on financial forecasting by statistical
average method,'' \emph{IEEE Access}, vol.~13, pp.~35355--35375, 2025.

\bibitem{choi2023} J. Choi, S. Yoo, X. Zhou and Y. Kim, ``Hybrid information mixing module
for stock movement prediction,'' \emph{IEEE Access}, vol.~11, pp.~28781--28790, 2023.

\bibitem{ho2021} T.-T. Ho and Y. Huang, ``Stock price movement prediction using sentiment
analysis and CandleStick chart representation,'' \emph{Sensors}, vol.~21, no.~23, p.~7957,
2021.

\bibitem{rouf2021} N. Rouf, M. B. Malik, T. Arif, S. Sharma, S. Singh, S. Aich and
H.-C. Kim, ``Stock market prediction using machine learning techniques: a decade survey on
methodologies, recent developments, and future directions,'' \emph{Electronics}, vol.~10,
no.~21, p.~2717, 2021.

\bibitem{song2024} Y. Song, ``Enhancing stock market trend reversal prediction using
feature-enriched neural networks,'' \emph{Heliyon}, vol.~10, p.~e24136, 2024.

\bibitem{javaid2024} N. Javaid, F. Aslam et al., ``Enhanced prediction of stock markets
using a novel deep learning model PLSTM-TAL in urbanized smart cities,'' \emph{Heliyon},
vol.~10, p.~e27747, 2024.

\bibitem{shaban2024} W. M. Shaban et al., ``SMP-DL: a novel stock market prediction
approach based on deep learning for effective trend forecasting,'' \emph{Neural Computing
and Applications}, vol.~36, pp.~1849--1873, 2024.

\bibitem{li2024llmsurvey} Y. Li, S. Wang, H. Ding and H. Chen, ``Large language models in
finance: a survey,'' 2024.

\bibitem{todd2024} A. Todd, J. Bowden and Y. Moshfeghi, ``Text-based sentiment analysis in
finance: synthesising the existing literature and exploring future directions,''
\emph{Intelligent Systems in Accounting, Finance and Management}, vol.~31, no.~1,
p.~e1549, 2024.

\bibitem{bahoo2024} S. Bahoo, M. Cucculelli, X. Goga and J. Mondolo, ``Artificial
intelligence in finance: a comprehensive review through bibliometric and content
analysis,'' \emph{SN Business and Economics}, vol.~4, no.~23, 2024.

\bibitem{khan2024chatgpt} M. S. Khan and H. Umer, ``ChatGPT in finance: applications,
challenges, and solutions,'' \emph{Heliyon}, vol.~10, no.~2, p.~e24890, 2024.

\bibitem{fatouros2024} G. Fatouros, J. Soldatos, K. Kouroumali, G. Makridis and
D. Kyriazis, ``Transforming sentiment analysis in the financial domain with ChatGPT,''
2024.

\bibitem{wilksch2024} M. Wilksch and O. Abramova, ``PyFin-sentiment: towards a
machine-learning-based model for deriving sentiment from financial tweets,'' 2024.

\bibitem{wu2023bloomberggpt} S. Wu, O. Irsoy, S. Lu, V. Dabravolski, M. Dredze,
S. Gehrmann, P. Kambadur, D. Rosenberg and G. Mann, ``BloombergGPT: a large language model
for finance,'' 2023.

\bibitem{yang2023fingpt} H. Yang, X.-Y. Liu and C. D. Wang, ``FinGPT: open-source
financial large language models,'' arXiv preprint, 2023.

\bibitem{wang2024fingptbench} N. Wang, H. Yang and C. D. Wang, ``FinGPT: instruction
tuning benchmark for open-source large language models in financial datasets,'' 2024.

\bibitem{araci2019} D. Araci, ``FinBERT: financial sentiment analysis with pre-trained
language models,'' arXiv:1908.10063, 2019.

\bibitem{hutto2014} C. J. Hutto and E. Gilbert, ``VADER: a parsimonious rule-based model
for sentiment analysis of social media text,'' in \emph{Proc. 8th Int. AAAI Conf. Weblogs
and Social Media}, 2014.

\bibitem{loughran2011} T. Loughran and B. McDonald, ``When is a liability not a liability?
Textual analysis, dictionaries, and 10-Ks,'' \emph{The Journal of Finance}, vol.~66,
no.~1, pp.~35--65, 2011.

\bibitem{vaswani2017} A. Vaswani, N. Shazeer, N. Parmar, J. Uszkoreit, L. Jones,
A. N. Gomez, {\L}. Kaiser and I. Polosukhin, ``Attention is all you need,'' in \emph{Adv.
Neural Inf. Process. Syst.}, 2017, pp.~5998--6008.

\bibitem{hochreiter1997} S. Hochreiter and J. Schmidhuber, ``Long short-term memory,''
\emph{Neural Computation}, vol.~9, no.~8, pp.~1735--1780, 1997.

\bibitem{chen2016xgboost} T. Chen and C. Guestrin, ``XGBoost: a scalable tree boosting
system,'' in \emph{Proc. 22nd ACM SIGKDD Int. Conf. Knowledge Discovery and Data Mining},
2016, pp.~785--794.

\bibitem{koa2024} K. J. L. Koa, Y. Ma, R. Ng and T.-S. Chua, ``Learning to generate
explainable stock predictions using self-reflective large language models,'' in
\emph{Proc. ACM Web Conference (WWW '24)}, Singapore, 2024.

\bibitem{lundberg2017} S. M. Lundberg and S.-I. Lee, ``A unified approach to interpreting
model predictions,'' in \emph{Adv. Neural Inf. Process. Syst.}, 2017, pp.~4765--4774.

\bibitem{ribeiro2016} M. T. Ribeiro, S. Singh and C. Guestrin, ```Why should I trust you?'
Explaining the predictions of any classifier,'' in \emph{Proc. 22nd ACM SIGKDD Int. Conf.
Knowledge Discovery and Data Mining}, 2016, pp.~1135--1144.

\bibitem{xu2026causal} S. Xu, Q. Cheng and C.-G. Lee, ``A causal perspective of stock
prediction models,'' \emph{IEEE Trans. Knowledge and Data Engineering}, vol.~38, no.~1,
pp.~14--25, 2026.

\bibitem{liu2026dynamic} H. Liu, B. Hu, Y. Zhou and Y. Zhou, ``A dynamic-causal lens of
stock interactions: graph modeling from symbolic movement patterns,'' \emph{IEEE Trans.
Knowledge and Data Engineering}, vol.~38, no.~2, pp.~722--735, 2026.

\bibitem{yu2025causalseqgnn} H. Yu and J. Liu, ``CausalSeqGNN: a Fama--French inspired
fuzzy cognitive map combined with graph neural network for stock prediction,'' \emph{IEEE
Trans. Fuzzy Systems}, vol.~33, no.~12, pp.~4418--4430, 2025.

\bibitem{luo2025magicnet} D. Luo, S. Li, W. Liao and R. Yan, ``MagicNet: memory-aware
graph interactive causal network for multivariate stock price movement prediction,''
\emph{IEEE Trans. Knowledge and Data Engineering}, vol.~37, no.~4, pp.~1989--2000, 2025.

\bibitem{wang2025mtmd} M. Wang, J. Tian, M. Zhang, J. Guo and W. Jia, ``MTMD: multi-scale
temporal memory learning and efficient debiasing framework for stock trend forecasting,''
\emph{IEEE Trans. Emerging Topics in Computational Intelligence}, vol.~9, no.~3,
pp.~2151--2163, 2025.

\bibitem{zhu2025multigran} P. Zhu, Y. Li, Q. Liu, D. Cheng and C. Jiang, ``Financial time
series prediction with multi-granularity graph augmented learning,'' \emph{IEEE Trans.
Knowledge and Data Engineering}, vol.~37, no.~11, pp.~6436--6449, 2025.

\bibitem{xing2024vague} R. Xing, R. Cheng, J. Huang, Q. Li and J. Zhao, ``Learning to
understand the vague graph for stock prediction with momentum spillovers,'' \emph{IEEE
Trans. Knowledge and Data Engineering}, vol.~36, no.~4, pp.~1698--1712, 2024.

\bibitem{patel2024hybrid} M. Patel, K. Jariwala and C. Chattopadhyay, ``A hybrid
relational approach toward stock price prediction and profitability,'' \emph{IEEE Trans.
Artificial Intelligence}, vol.~5, no.~11, pp.~5844--5854, 2024.

\bibitem{liu2024heterogeneous} H. Liu, Y. Zhou, Y. Zhou and B. Hu, ``Heterogeneous
dual-dynamic attention network for modeling mutual interplay of stocks,'' \emph{IEEE
Trans. Artificial Intelligence}, vol.~5, no.~7, pp.~3595--3606, 2024.

\bibitem{gao2025relational} Q. Gao, Z. Liu, L. Huang, K. Zhang, J. Wang and G. Liu,
``Relational stock selection via probabilistic state space learning,'' \emph{IEEE Trans.
Knowledge and Data Engineering}, vol.~37, no.~2, pp.~865--880, 2025.

\bibitem{kumar2024rdmdea} A. Kumar and S. Yadav, ``A multiobjective multiperiod efficient
portfolio selection approach for positive and negative returns using RDM-DEA under a
credibilistic framework,'' \emph{IEEE Trans. Engineering Management}, vol.~71,
pp.~14114--14125, 2024.

\bibitem{gu2025ragic} J. Gu, W. Du and G. Wang, ``RAGIC: risk-aware generative framework
for stock interval construction,'' \emph{IEEE Trans. Knowledge and Data Engineering},
vol.~37, no.~4, pp.~2085--2096, 2025.

\bibitem{zhu2025spin} M. Zhu et al., ``SPIN: sparse portfolio strategy with irregular news
in fluctuating markets,'' \emph{IEEE Trans. Knowledge and Data Engineering}, vol.~37,
no.~6, pp.~3714--3727, 2025.

\bibitem{elahi2024} A. Elahi and F. Taghvaei, ``Combining financial data and news articles
for stock price movement prediction using large language models,'' in \emph{Proc. IEEE
Int. Conf. Big Data}, 2024, pp.~4875--4883.

\bibitem{mishra2025} N. Mishra, M. Anzar, S. Pandey and S. Mishra, ``A review on stock
market trends and stock price prediction using sentiment analysis and market data,'' in
\emph{Proc. Int. Conf. Communication, Security, and Artificial Intelligence}, 2025,
pp.~56--62.

\bibitem{officioso2026} G. Officioso, G. Sperl\`i and A. Vignali, ``Financial news
sentiment meets market data: a large language model-based approach to stock price
prediction,'' \emph{Information Sciences}, vol.~754, p.~123711, 2026.

\bibitem{zong2025} C. Zong, J. Wan, L. Cascone and H. Zhou, ``Stock movement prediction
with multimodal stable fusion via gated cross-attention mechanism,'' \emph{Complex and
Intelligent Systems}, vol.~11, no.~9, 2025.

\bibitem{singsiri2026} P. Singsiri and J. Yokrattanasak, ``Machine learning-driven
portfolio optimization using money flow index-based sentiment signals,''
\emph{International Journal of Financial Studies}, vol.~14, no.~5, p.~112, 2026.

\bibitem{naidoo2025} T. Naidoo, P. Moores-Pitt, P.-F. Muzindutsi and K. O. Isah,
``Analysing investor sentiment and stock market volatility of the JSE size-based indices:
a GARCH-MIDAS approach,'' \emph{Risk Management}, vol.~27, no.~3, 2025.

\bibitem{markowitz1952} H. Markowitz, ``Portfolio selection,'' \emph{The Journal of
Finance}, vol.~7, no.~1, pp.~77--91, 1952.

\bibitem{ledoit2004} O. Ledoit and M. Wolf, ``A well-conditioned estimator for
large-dimensional covariance matrices,'' \emph{Journal of Multivariate Analysis}, vol.~88,
no.~2, pp.~365--411, 2004.

\bibitem{demiguel2009} V. DeMiguel, L. Garlappi and R. Uppal, ``Optimal versus naive
diversification: how inefficient is the $1/N$ portfolio strategy?'' \emph{The Review of
Financial Studies}, vol.~22, no.~5, pp.~1915--1953, 2009.

\bibitem{lopezdeprado2016} M. L\'opez de Prado, ``Building diversified portfolios that
outperform out of sample,'' \emph{The Journal of Portfolio Management}, vol.~42, no.~4,
pp.~59--69, 2016.

\bibitem{hamilton1989} J. D. Hamilton, ``A new approach to the economic analysis of
nonstationary time series and the business cycle,'' \emph{Econometrica}, vol.~57, no.~2,
pp.~357--384, 1989.

\bibitem{ang2002} A. Ang and G. Bekaert, ``International asset allocation with regime
shifts,'' \emph{The Review of Financial Studies}, vol.~15, no.~4, pp.~1137--1187, 2002.

\bibitem{kritzman2012} M. Kritzman, S. Page and D. Turkington, ``Regime shifts:
implications for dynamic strategies,'' \emph{Financial Analysts Journal}, vol.~68, no.~3,
pp.~22--39, 2012.

\bibitem{luo2026regime} Y. Luo and J. M. Mulvey, ``Regime-aware asset allocation with
dual-regime signals and regime-dependent asset selection,'' \emph{The Journal of Financial
Data Science}, vol.~8, no.~3, pp.~104--134, 2026.

\bibitem{moreira2017} A. Moreira and T. Muir, ``Volatility-managed portfolios,'' \emph{The
Journal of Finance}, vol.~72, no.~4, pp.~1611--1644, 2017.

\bibitem{bakerwurgler2007} M. Baker and J. Wurgler, ``Investor sentiment in the stock
market,'' \emph{Journal of Economic Perspectives}, vol.~21, no.~2, pp.~129--151, 2007.

\bibitem{grinold2000} R. C. Grinold and R. N. Kahn, \emph{Active Portfolio Management},
2nd~ed. New York: McGraw-Hill, 2000.

\bibitem{gu2020} S. Gu, B. Kelly and D. Xiu, ``Empirical asset pricing via machine
learning,'' \emph{The Review of Financial Studies}, vol.~33, no.~5, pp.~2223--2273, 2020.

\bibitem{harvey2015} C. R. Harvey and Y. Liu, ``Backtesting,'' \emph{The Journal of
Portfolio Management}, vol.~42, no.~1, pp.~13--28, 2015.

\bibitem{valls2026} F. B. Valls and E. Steurer, ``When strategies work too well:
volatility-based investing, overfitting and the limits of empirical finance,''
\emph{Journal of Financial Risk Management}, vol.~15, no.~3, pp.~189--211, 2026.

\bibitem{newey1987} W. K. Newey and K. D. West, ``A simple, positive semi-definite,
heteroskedasticity and autocorrelation consistent covariance matrix,''
\emph{Econometrica}, vol.~55, no.~3, pp.~703--708, 1987.

\bibitem{politis1994} D. N. Politis and J. P. Romano, ``The stationary bootstrap,''
\emph{Journal of the American Statistical Association}, vol.~89, no.~428,
pp.~1303--1313, 1994.

\bibitem{granger1969} C. W. J. Granger, ``Investigating causal relations by econometric
models and cross-spectral methods,'' \emph{Econometrica}, vol.~37, no.~3, pp.~424--438,
1969.

\bibitem{spirtes2000} P. Spirtes, C. Glymour and R. Scheines, \emph{Causation,
Prediction, and Search}, 2nd~ed. Cambridge, MA: MIT Press, 2000.

\end{thebibliography}
\end{multicols}
\end{scriptsize}
